\chapter{哈佛教授解釋寫作的規則}

這場講座一開始便提出一個實際的承諾。我們並不是被要求為了語言理論本身而進入理論；我們被問的是：如何把文章寫好，以及在大型語言模型已經能寫出流暢、稱職文句的時代，如何持續把文章寫好。Pinker 回答這個問題時，並不是丟出一袋零碎技巧，而是提出一連串因果主張、對照與完整範例。這場講座沒有真正的黑板數學，因此下文中的公式都是根據講座中以散文表述的機制所做的示意性重建。它們的目的，是在忠於逐字稿的前提下，讓其中的邏輯更清晰。

\section{糟糕寫作與知識的詛咒}

訪談者一開始便直指那個主導性問題：為什麼糟糕的寫作如此之多？Pinker 先清除掉一個最討喜、也最容易讓人自鳴得意的錯誤答案。人們很容易說，不透明的散文主要是一種策略：作者用術語掩飾單薄的想法，以此顯得深奧、保護自己的小圈子，或把外人排除在外。Pinker 並不否認這種情況確實會發生。但他堅持，這並不是主要解釋。他認識太多真正聰明、明明有很多可說之物的人，卻仍然寫出令人難以卒讀的文章。

這個反轉很重要，因為它讓整場講座踏上正確軌道。核心的失敗，不是一場道德戲劇，而是一種校準失準。

\begin{definition}
所謂知識的詛咒，是指我們難以重新構想：對於那些並不知道我們所知之事的人而言，世界看起來究竟是什麼樣子。
\end{definition}

Pinker 把這種失敗濃縮為一個作者本應自問、卻往往來不及問清楚的問題：讀者知道什麼？又不知道什麼？若把它寫成一個示意性的機制，我們便得到
\begin{equation}
\text{專家知識} \to \text{無法為讀者建模} \to \text{術語／縮寫／抽象語} \to \text{讀者困惑}.
\end{equation}

\begin{remark}
這條公式是編者的速記，而不是講座中可見的記號。它的目的，只是把 Pinker 的因果鏈條明白攤開。
\end{remark}

接著，講座以一個例子讓這條公式站住腳。某位傑出的分子生物學家做 TED 演講時，說話的方式完全像是在對其他分子生物學家發言。他並不從問題講起，也不解釋這項工作為何重要，而是直接一頭栽進實驗細節裡。台下幾乎立刻就跟不上了，而唯一沒看見這一點的人，正是講者自己。Pinker 的意思既嚴厲又準確：這個人並不愚蠢，除了在一個局部、卻災難性的地方。他不知道，對於那些不知道他所知之事的人而言，世界是什麼樣子。

同樣的失敗，也出現在抽象語快過解釋的時候。Pinker 引了一句話：``the level of the stimulus was proportional to the intensity of the reaction.'' 在他舉的那個例子裡，這句話實際上只是說：孩子看小兔子的時間，比看卡車的時間更久。前一句未必不真，但它把場景藏起來了。後一句則給了讀者一個可以看見的東西。

\subsection{Question \& Answer}

\paragraph{Question.} 如果許多作者都很聰明，也很真誠，為什麼糟糕的寫作還是這麼多？

\paragraph{Answer.} 因為在某個領域裡的聰明，並不保證你有能力定位讀者。作者的知識跑得比作者對受眾的模型更快。Pinker 最核心的動作，就是用一個關於視角代入失敗的故事，取代那個關於虛榮或惡意的故事。

\begin{proposition}
在解釋性散文中，最主要的失敗模式並不是缺乏想法，而是無法估計讀者究竟能合理推知到什麼程度。
\end{proposition}

這也就是為什麼 Pinker 會把知識的詛咒連到自我中心傾向，以及一種缺席的心智理論。於是，講座一開始便由一種認知錯誤切入，而這個錯誤將支配後面每一個關於例子、結構、節奏與風格的問題。

\section{把受眾校準當作一種寫作實踐}

一旦診斷成立，講座立刻追問補救之道。若糟糕的寫作是受眾模型失效，那麼我們在起草時究竟該做什麼？Pinker 的第一個回答很克制：試著想像讀者，培養同理心。但他幾乎立刻又補充說，這還不夠。知識的詛咒之所以狡猾，正因為從內部感受它時，你是感覺不到它的。那些你以為太顯然、無須解釋的東西，往往其實一點也不顯然。

若用示意方式寫，修改過程大致如下：
\begin{align}
\text{草稿} &\to \text{想像中的讀者}, \\
\text{草稿} + \text{真實讀者的反應} &\to \text{修改}.
\end{align}
第一條箭頭不可或缺；第二條箭頭才真正把我們救出來。

Pinker 在這裡特別高明，因為他拒絕把讀者想成一個粗糙的形象。他當年並不是因為母親遲鈍或缺乏修養，才把草稿給她看；他說的恰恰相反。她很聰明，讀書很多，也很認真。真正重要的是，她不是他所屬領域的專家。合適的測試讀者，不是隨便抓來的人，也不是同行，而是一個足夠聰明，一旦論點被妥善表述便能跟上的人；但同時又離得夠遠，不共享作者那套局部速記系統。

講座於是走過一串非常具體的讀者序列：
\begin{itemize}
\item 作者自己所做的同理心嘗試，
\item 一位值得信任、聰明的非專業者，
\item 商業出版社的編輯，
\item 鄰近領域的學者，
\item 甚至是那些剛好站在作者極小研究圈外的學生或同事。
\end{itemize}

最後這一點既重要，又很容易被忽略。Pinker 指出，即便在學院內部，有時甚至就在同一個系裡，人們也會變得彼此不可理解，因為實驗室裡五六個人已經一起呼吸同一套術語好幾個月了。一旦這篇文字離開那個微型氣候，它立刻就會崩塌。

因此，講座的第一條實用規則，不僅僅是 ``了解你的受眾''，而是更嚴格的：
\begin{quote}
試著進入讀者的頭腦，但不要只憑這種努力就相信自己已經成功。把草稿交到一個真正的人手上：那個人要聰明、好奇，並且不在你的直接小圈子裡。
\end{quote}

這是第一個主要的修復機制。第二個也隨之自然到來。

\section{具體性、意象與感官式理解}

講座在這裡轉了一個對整章都很重要的彎。一旦 Pinker 談完受眾校準，他便追問：理解本身究竟由什麼構成？他的回答是：語言固然重要，但它的重要性常常被高估。我們作為作者，棲身於語言之中，Pinker 自己更是專業地研究語言；然而，理解並不只是接收一串又一串的詞句。語言是用來讓讀者掌握某個想法的工具，而這個想法常常是視覺的、動作性的、身體性的、情感性的、聽覺的，或以其他感官形式存在的。

我們可以把這個主張壓縮為
\begin{equation}
\text{語言} \to \text{心理模型},
\qquad
\text{而不僅僅是}
\qquad
\text{語言} \to \text{更多語言}.
\end{equation}

這便是 Pinker 的第二條主要規則。如果第一條規則是找到讀者的位置，那麼第二條規則就是：給讀者某種可以搭建的東西。若你真正想說的是 ``兔子''，那就不要說 ``stimulus''。當讀者需要的是一個場景時，就不要躲在 level、perspective、framework、paradigm 或 concept 這些字後面。

\paragraph{A worked example.}
講座提供了一個幾乎完美的小型推演，示範如何從抽象走向可觀察性：
\begin{enumerate}
\item 先從這句話開始：``the level of the stimulus was proportional to the intensity of the reaction.''
\item 注意到，這句話記錄了一種關係，卻壓抑了幾乎所有可感知的內容。
\item 追問：讀者究竟應該在腦中看見什麼？
\item 把具體實體與動作補回來：孩子、兔子、卡車、看得比較久。
\item 於是，這個主張在散文中變得可見，因此也變得可思。
\end{enumerate}

若示意地寫，便是
\begin{align}
\text{``The level of the stimulus was proportional to the intensity of the reaction.''}
&\not\Rightarrow \text{給讀者一個穩定的場景}, \\
\text{``Kids look longer at a bunny than at a truck.''}
&\Rightarrow \text{一幅可供檢視的心理圖像}.
\end{align}

講座並沒有停在這一個例子上，而是繼續向外推到視覺隱喻。Pinker 認為，早期的散文之所以常給我們一種更鮮明的感覺，其中一個原因是：早期作者手上可供直接取用的現成抽象語比較少，因此不得不倚賴大家共同持有的圖像。他舉出 ``the spirit of the hawk kneaded into our flesh'' 這樣一個強烈的片語，來說明一種表達方式：它所做的工作，在後來可能已被一個像 ``aggression'' 或 ``antisocial behavior'' 這樣的技術術語所接管。我們不必整體地模仿舊風格，但應當理解其機制：作者是透過呈現某種可想像之物，來抵達讀者的。

\section{為什麼寫作比說話更難}

到了這裡，講座引入一個新的障礙。既然說話對兒童而言如此自然，為什麼寫作卻始終緩慢、費力，甚至帶著人工性？Pinker 的回答，是把層次往上移。到目前為止，他給的是建議；而在這裡，他開始解釋：正是在什麼樣的環境中，這些建議才顯得必要。

\subsection{Question \& Answer}

\paragraph{Question.} 為什麼寫作相較於說話，總讓人感到不自然？

\paragraph{Answer.} 因為對話發生於一個共享情境之中，而寫作卻必須用頁面上的符號，自己建構出那個情境。

講座中的對照可以再次示意地寫成
\begin{align}
\text{對話} &= \text{共同背景} + \text{指示語} + \text{回饋}, \\
\text{寫作} &= \text{頁面上的文字}.
\end{align}

這一小對等式，把逐字稿中幾個不同的要點濃縮了起來：
\begin{enumerate}
\item 對話不會無中生有地開始。參與者知道自己為何在場，也知道是什麼情境把他們帶到一起。
\item 由於情境是共享的，像 ``這個''、``那個''、``她''、``那件事''、``我剛才說的那個'' 這類詞，常常都完全清楚。
\item 對話還提供即時修正：皺起的眉頭、困惑的眼神、要求澄清的提問、注意力飄移的跡象。
\item 即使是現場公開演講，仍保留了部分這種支撐，因為講者還能看著聽眾。
\item 寫作則把這一切支撐拿掉。讀者可能住在別處、稍後才讀，對文字寫成時的在地情境一無所知。
\end{enumerate}

講座的節奏就建立在這一拍上。若少了它，例子與明確框架聽起來就會像是可有可無、甚至只是裝飾。有了它，它們便變成結構性的必需。寫作之所以比說話更難，是因為紙頁必須獨自承擔那些在對話中本由情境、手勢與回饋所承擔的負荷。

\section{概括、例子與意義的形狀}

現在，講座從大層次的對比收窄到一種我們可以逐句使用的方法。訪談者以一種刻意尖銳的方式提出一條原則：沒有例子的概括，以及沒有概括的例子，兩者都沒有用。Pinker 稍稍放軟了這個判詞，但接受了它的結構。沒有案例的概括，常常讓讀者點頭，卻其實並不知道究竟被主張了什麼。沒有概括的例子，則會讓讀者納悶：這個例子到底為什麼會被引進來？

我們可以把這個平衡概括為
\begin{equation}
\text{例子} + \text{概括} \to \text{理解}.
\end{equation}
接著，訪談者又把同一種平衡改述為 context 與 compression 之間的關係。這雖然不是 Pinker 自己的術語，但作為對這段交換的編者式摘要，卻相當有用：
\begin{equation}
\text{脈絡} \leftrightarrow \text{壓縮}.
\end{equation}

其中的邏輯按步展開：
\begin{enumerate}
\item 概括是一種壓縮；它跨越具體細節。
\item 正因它壓抑了細節，它可能無法喚起任何確定的所指。
\item 例子則重新給出一個具體範圍。
\item 但例子本身仍未告訴我們，究竟學到了什麼。
\item 概括於是回來，告訴我們這個例子為何重要。
\end{enumerate}

在這裡，講座做的事情比一般的寫作課更細膩。它從說明轉向語義學。Pinker 舉出一連串熟悉的表達，其意義並不能只靠把構成部分的意思加總起來而恢復。Bathroom 不一定有 bath。Going to the bathroom 也不是真正字面上地走向一間有浴缸的房間。Breakfast 未必會被體驗成「打破禁食」。Christmas 這個詞每次被說出來時，人們通常也不會重新把它當成一塊新鮮的組合神學來處理。

接著，講座又再次展開。Olive oil 是由橄欖做成的油，而 baby oil 卻是給嬰兒用的油。同樣的表面形式，支持著不同的語義關係。Richard Lederer 更帶玩味的例子，則把這個教訓推得更遠。Singer 清楚反映出一條仍然活著的規則：``-er'' 可以把動詞變成典型施事者的名詞。Finger 則不是。語言會保留古老的構詞形式，而它們原初的邏輯，卻早已從活的意識中消失。

\begin{remark}
意義往往是約定俗成的，也是歷史性的，而不是機械式地由成分組合而成。Pinker 的意思並不是說完全沒有規則，而是說，讀者不能被期待僅憑表面形式就恢復出作者想表達的意思。
\end{remark}

講座確實短暫提到 adultery 與 adulterate，用來說明語義漂移。對於這一刻，最審慎的筆記方式，就是不要讓這個例子承載超出講座本身所賦予它的重量。重要的只是：熟悉的字詞可以離原初關係漂移得很遠，而作者絕不應假定，單靠形式拆解就足以產生理解。

\section{風格作為愉悅：節奏、聲音、簡潔與幽默}

一旦講座已經顯示出：意義並不會自動由文字承載，它便準備好轉入風格。Pinker 對次序很小心。他並不把愉悅當成清晰的替代物，而是把它當作在清晰已經獲得保障之後的下一層。

在講座進入 Shakespeare 與 Strunk 之前，它先停留在兒童解釋中的新鮮感上。Pinker 很喜歡像 ``smoke is fire vapor'' 這樣的例子，因為它們顯示的是一顆仍在伸手抓取可見關係、而非現成陳腔的心靈。孩子尚未擁有那一整套後來作者常用來避難的厚重抽象庫。從這個意義上說，他們有點像早期的作者：那時人們不得不訴諸共同圖像，因為還不能倚賴後來那種技術詞彙。這堂課的教訓不是叫我們像孩子一樣寫，而是要我們明白：新鮮感與可觀察性之間，關係極深。

\subsection{Question \& Answer}

\paragraph{Question.} 什麼能讓散文不只是可理解，而且令人愉悅？

\paragraph{Answer.} Pinker 的答案是：散文有一種可聽見的生命。它可以走得好，也可以走得壞；它可以帶著節奏，也可以把節奏弄斷；它可以刺耳，也可以滑順。作者可以直接測試這一點：把文章朗讀出來，聽出摩擦所在，再修改那一行，直到它能被乾淨地說出來。

局部的更新規則很簡單：
\begin{equation}
\text{朗讀} \to \text{聽見摩擦} \to \text{修改節奏}.
\end{equation}

於是，若干風格變數現在作為操作性量，而不是空泛讚語，被引入場中：
\begin{itemize}
\item \textbf{視覺意象}：這個句子能否喚起一個場景？
\item \textbf{韻律結構}：散文是否以可用的節拍前進，而不是步履蹣跚？
\item \textbf{噝音}：是否連續出現太多帶噝聲的音？
\item \textbf{頭韻}：是否有一點輕微的聲音重現，帶來愉悅的火花，卻又不至於用力過猛？
\item \textbf{簡潔}：同樣的意思，能否用更少的拖沓傳達？
\end{itemize}

Pinker 也把幽默直接接到這一節上。幽默依賴新鮮感，也依賴時機。笑話一旦拖得太長，就不再好笑；包袱若被預先大張旗鼓地鋪陳，也會在自己的準備之下坍塌。因此，《Hamlet》裡那句 ``brevity is the soul of wit'' 一語兩用：它既陳述了關於機智的真理，也以自身的經濟性示範了那條真理。

Strunk 的 ``omit needless words'' 也是以同樣方式被處理。它不只是一本手冊裡隨手摘來的標語。它本身就是這條原則的例子，而 Pinker 又補上一個來自自己寫作生活中的實際檢驗：當報紙或雜誌施加嚴格的字數上限時，被迫做出的壓縮，往往不是單純縮短文章，而是讓文章變得更好。

若示意地寫，
\begin{equation}
\text{不必要的字詞} \to \text{額外處理負擔} \to \text{力道減弱}.
\end{equation}

講座在這裡的論證，既是認知性的，也是美學性的。每多一個字，都是讀者多做一分工。除此之外，壓縮往往還會讓節奏更銳利，迫使作者離開那些綿軟含糊的慣用語，而走向更具體的措辭。依這種觀點，風格是一種有紀律的愉悅。

\section{舊散文、學術散文與 LLM 散文}

講座此時從技藝向外擴展到制度、歷史與科技。Pinker 先解釋，為什麼他如此在意糟糕的學術散文。這種抱怨不只是風格潔癖。糟糕的散文會浪費智力，也會浪費公共價值。它迫使審稿者、同事與學生，把時間耗在解碼那些本來就應該在源頭處寫清楚的句子上。一段必須讀上五六遍，主張才終於浮出水面的文字，不只是醜陋而已；它既低效，又危險。

這種制度性的抱怨，接著又打開一個歷史性的問題。為什麼古老的散文如此常讓現代讀者覺得鮮明，甚至繁麗？Pinker 依序提出幾個原因。早期作者受的是強而有力的文學典範教育。散文更公開地被培養成一種自我呈現於世界的媒介。而對本講座內在邏輯而言最重要的一點是：早期作者手邊現成的抽象語與陳腔較少，因此必須把新想法鑄造成普通讀者也能看見的形式。

接著，這個歷史論點又轉向文化變遷。Pinker 提到一個漫長的 informalization 過程：階級感更少、儀式性更弱、對話式親密更多，對任何聽起來高昂或刻意打磨過的語調也更容易起疑。這種轉變有助於解釋：為什麼許多現代作者即使完全有能力寫出強而有工藝感的散文，卻仍然猶疑，不敢讓自己的文字聽起來太正式、太完成。

直到這一切鋪墊之後，講座才回到 AI。這個順序很重要。LLM 散文不是作為某種外來異物被引進，而是作為一種施加在既有問題上的最新壓力被引進。Pinker 先承認它的長處。它往往是清楚的，至少在有限意義上如此：它的句子有序、句法平易、段落通常也有可辨認的開頭與結尾。但它的弱點同樣醒目：它很泛化、很散文化，而且常常讓人一眼就看出那是某種總體風格的仿作。

我們可以把這種對照示意地寫成
\begin{align}
\text{LLM 散文} &= \text{清楚的排序} + \text{平易的句法} + \text{泛化的仿作}, \\
\text{強而有力的散文} &= \text{清晰} + \text{具體} + \text{新鮮}.
\end{align}

因此，講座最後落在一條重要的不等價關係上：
\begin{equation}
\text{清楚的散文} \neq \text{新鮮的散文}.
\end{equation}

Pinker 也短暫推測，為什麼 LLM 散文往往比學者、律師或官僚的文字更容易讀。也許是因為它在訓練與回饋中被反覆捶打成形。也許，更帶推測性地說，大量例子的綜合消除了若干最糟糕的局部扭曲。但他並沒有讓這種推測飄向超出講座所能支持的更強主張。他沒有下結論說，人類心智不過就是一個大型語言模型。相反地，他抽出一個更狹窄的教訓：任何嚴肅的智能理論，如今都必須比那些較舊、偏重規則的觀點，更重視從龐大輸入中提取大尺度模式的能力。

\section{Summary}

整場講座是以一條鏈的方式展開的，而最好也這樣來讀。它一開始便拒絕那個令人自滿的故事：糟糕散文主要是策略性的晦澀。它用知識的詛咒取而代之，接著立刻談修復：真實讀者、真實編輯、真實測試。然後，它把論證推得更深，追問理解究竟是什麼，並回答說：讀者需要的不只是文字，而是心理模型。這個回答進一步引向說話與寫作之間的不對稱，再從那裡走向例子與概括之間的平衡。一旦講座已經表明，意義既不是自動的，也不是完全可由組合規則推出的，它就能進入風格，把風格當作一種可操作的東西：節奏、聲音、壓縮、機智與愉悅。最後，講座再把視野擴展到歷史、制度與 AI，同時又始終不失起初的那個核心。從第一頁到最後一頁，作者的任務始終如一：把我們真正想說的東西，說成另一顆心智可以真實進入的形式。